% bg-template.tex — Community pattern template for the Wildcat Beamer theme
%
% HOW TO USE
% ----------
% Add this line to your preamble, AFTER \usetheme{wildcat} (adjust the path
% and filename as needed):
%
%   % bg-template.tex — Community pattern template for the Wildcat Beamer theme
%
% HOW TO USE
% ----------
% Add this line to your preamble, AFTER \usetheme{wildcat} (adjust the path
% and filename as needed):
%
%   % bg-template.tex — Community pattern template for the Wildcat Beamer theme
%
% HOW TO USE
% ----------
% Add this line to your preamble, AFTER \usetheme{wildcat} (adjust the path
% and filename as needed):
%
%   % bg-template.tex — Community pattern template for the Wildcat Beamer theme
%
% HOW TO USE
% ----------
% Add this line to your preamble, AFTER \usetheme{wildcat} (adjust the path
% and filename as needed):
%
%   \input{bg-gallery/bg-template} 
%
% Note: use \input{}, not \include{}. The \include command only works inside
% the document body; pattern commands must be set in the preamble so that
% the theme's background saveboxes are populated with the correct pattern.
%
%
% HOW TO CONTRIBUTE A PATTERN
% ----------------------------
% 1. Copy this file and give it a descriptive name (e.g. bg-chevron.tex).
% 2. Rename \bgpatterntemplate to something unique (e.g. \bgpatternchevron).
% 3. Replace the TikZ drawing commands in \bgpatterntemplate with your own.
%    The command receives one argument (#1) — the color — so the theme can
%    reuse it in different contexts (title, section slides, tcolorbox blocks).
%    Your drawing should fill the full slide; the enclosing tikzpicture
%    bounding box is already set to (0,0)--(\paperwidth,\paperheight).
% 4. Submit a pull request to https://github.com/aarondwolf/wildcat
%
%
% PATTERN COMMAND REQUIREMENTS
% -----------------------------
% \bgpatterntemplate{color}  — draws a full-slide background given a color name.
%   - Must work inside an existing tikzpicture (no \begin{tikzpicture} wrapper).
%   - The bounding box already covers the full slide; just draw your shapes.
%   - Use #1 as the color argument wherever you reference the theme color.
%   - Lighter/darker shades (#1!30!white, #1!80!black, etc.) are fine to use.
%
% -----------------------------------------------------------------------
% EXAMPLE PATTERN: plain solid fill
% -----------------------------------------------------------------------
% This is the simplest possible pattern — a single solid rectangle covering
% the entire slide. Replace the \fill command below with your own TikZ
% drawing to create a new pattern.

\newcommand{\bgpatterntemplate}[1]{%
  \fill[#1] (0,0) rectangle (\paperwidth,\paperheight);%
  % Add your Tikz commands here. For color mising, use
  % the pattern #1!N!<colorname>, where N is a number
  % from 0 to 100 indicating the percentage of mixing. 
  % For example:
  %   #1!30!white — 30% white, 70% base color
  %   #1!80!black — 80% black, 20% base color
}

% Wire the pattern into the theme. Both lines are required:
%   \bgpattern      — used for the title page, section slides, and standout frames
%   \bgpatternblock — used for pblock/palertblock/pexampleblock title bars
%                     (receives the block color as its argument)
\renewcommand{\bgpattern}{\bgpatterntemplate{wcprimary}}
\renewcommand{\bgpatternblock}[1]{\bgpatterntemplate{#1}}
 
%
% Note: use \input{}, not \include{}. The \include command only works inside
% the document body; pattern commands must be set in the preamble so that
% the theme's background saveboxes are populated with the correct pattern.
%
%
% HOW TO CONTRIBUTE A PATTERN
% ----------------------------
% 1. Copy this file and give it a descriptive name (e.g. bg-chevron.tex).
% 2. Rename \bgpatterntemplate to something unique (e.g. \bgpatternchevron).
% 3. Replace the TikZ drawing commands in \bgpatterntemplate with your own.
%    The command receives one argument (#1) — the color — so the theme can
%    reuse it in different contexts (title, section slides, tcolorbox blocks).
%    Your drawing should fill the full slide; the enclosing tikzpicture
%    bounding box is already set to (0,0)--(\paperwidth,\paperheight).
% 4. Submit a pull request to https://github.com/aarondwolf/wildcat
%
%
% PATTERN COMMAND REQUIREMENTS
% -----------------------------
% \bgpatterntemplate{color}  — draws a full-slide background given a color name.
%   - Must work inside an existing tikzpicture (no \begin{tikzpicture} wrapper).
%   - The bounding box already covers the full slide; just draw your shapes.
%   - Use #1 as the color argument wherever you reference the theme color.
%   - Lighter/darker shades (#1!30!white, #1!80!black, etc.) are fine to use.
%
% -----------------------------------------------------------------------
% EXAMPLE PATTERN: plain solid fill
% -----------------------------------------------------------------------
% This is the simplest possible pattern — a single solid rectangle covering
% the entire slide. Replace the \fill command below with your own TikZ
% drawing to create a new pattern.

\newcommand{\bgpatterntemplate}[1]{%
  \fill[#1] (0,0) rectangle (\paperwidth,\paperheight);%
  % Add your Tikz commands here. For color mising, use
  % the pattern #1!N!<colorname>, where N is a number
  % from 0 to 100 indicating the percentage of mixing. 
  % For example:
  %   #1!30!white — 30% white, 70% base color
  %   #1!80!black — 80% black, 20% base color
}

% Wire the pattern into the theme. Both lines are required:
%   \bgpattern      — used for the title page, section slides, and standout frames
%   \bgpatternblock — used for pblock/palertblock/pexampleblock title bars
%                     (receives the block color as its argument)
\renewcommand{\bgpattern}{\bgpatterntemplate{wcprimary}}
\renewcommand{\bgpatternblock}[1]{\bgpatterntemplate{#1}}
 
%
% Note: use \input{}, not \include{}. The \include command only works inside
% the document body; pattern commands must be set in the preamble so that
% the theme's background saveboxes are populated with the correct pattern.
%
%
% HOW TO CONTRIBUTE A PATTERN
% ----------------------------
% 1. Copy this file and give it a descriptive name (e.g. bg-chevron.tex).
% 2. Rename \bgpatterntemplate to something unique (e.g. \bgpatternchevron).
% 3. Replace the TikZ drawing commands in \bgpatterntemplate with your own.
%    The command receives one argument (#1) — the color — so the theme can
%    reuse it in different contexts (title, section slides, tcolorbox blocks).
%    Your drawing should fill the full slide; the enclosing tikzpicture
%    bounding box is already set to (0,0)--(\paperwidth,\paperheight).
% 4. Submit a pull request to https://github.com/aarondwolf/wildcat
%
%
% PATTERN COMMAND REQUIREMENTS
% -----------------------------
% \bgpatterntemplate{color}  — draws a full-slide background given a color name.
%   - Must work inside an existing tikzpicture (no \begin{tikzpicture} wrapper).
%   - The bounding box already covers the full slide; just draw your shapes.
%   - Use #1 as the color argument wherever you reference the theme color.
%   - Lighter/darker shades (#1!30!white, #1!80!black, etc.) are fine to use.
%
% -----------------------------------------------------------------------
% EXAMPLE PATTERN: plain solid fill
% -----------------------------------------------------------------------
% This is the simplest possible pattern — a single solid rectangle covering
% the entire slide. Replace the \fill command below with your own TikZ
% drawing to create a new pattern.

\newcommand{\bgpatterntemplate}[1]{%
  \fill[#1] (0,0) rectangle (\paperwidth,\paperheight);%
  % Add your Tikz commands here. For color mising, use
  % the pattern #1!N!<colorname>, where N is a number
  % from 0 to 100 indicating the percentage of mixing. 
  % For example:
  %   #1!30!white — 30% white, 70% base color
  %   #1!80!black — 80% black, 20% base color
}

% Wire the pattern into the theme. Both lines are required:
%   \bgpattern      — used for the title page, section slides, and standout frames
%   \bgpatternblock — used for pblock/palertblock/pexampleblock title bars
%                     (receives the block color as its argument)
\renewcommand{\bgpattern}{\bgpatterntemplate{wcprimary}}
\renewcommand{\bgpatternblock}[1]{\bgpatterntemplate{#1}}
 
%
% Note: use \input{}, not \include{}. The \include command only works inside
% the document body; pattern commands must be set in the preamble so that
% the theme's background saveboxes are populated with the correct pattern.
%
%
% HOW TO CONTRIBUTE A PATTERN
% ----------------------------
% 1. Copy this file and give it a descriptive name (e.g. bg-chevron.tex).
% 2. Rename \bgpatterntemplate to something unique (e.g. \bgpatternchevron).
% 3. Replace the TikZ drawing commands in \bgpatterntemplate with your own.
%    The command receives one argument (#1) — the color — so the theme can
%    reuse it in different contexts (title, section slides, tcolorbox blocks).
%    Your drawing should fill the full slide; the enclosing tikzpicture
%    bounding box is already set to (0,0)--(\paperwidth,\paperheight).
% 4. Submit a pull request to https://github.com/aarondwolf/wildcat
%
%
% PATTERN COMMAND REQUIREMENTS
% -----------------------------
% \bgpatterntemplate{color}  — draws a full-slide background given a color name.
%   - Must work inside an existing tikzpicture (no \begin{tikzpicture} wrapper).
%   - The bounding box already covers the full slide; just draw your shapes.
%   - Use #1 as the color argument wherever you reference the theme color.
%   - Lighter/darker shades (#1!30!white, #1!80!black, etc.) are fine to use.
%
% -----------------------------------------------------------------------
% EXAMPLE PATTERN: plain solid fill
% -----------------------------------------------------------------------
% This is the simplest possible pattern — a single solid rectangle covering
% the entire slide. Replace the \fill command below with your own TikZ
% drawing to create a new pattern.

\newcommand{\bgpatterntemplate}[1]{%
  \fill[#1] (0,0) rectangle (\paperwidth,\paperheight);%
  % Add your Tikz commands here. For color mising, use
  % the pattern #1!N!<colorname>, where N is a number
  % from 0 to 100 indicating the percentage of mixing. 
  % For example:
  %   #1!30!white — 30% white, 70% base color
  %   #1!80!black — 80% black, 20% base color
}

% Wire the pattern into the theme. Both lines are required:
%   \bgpattern      — used for the title page, section slides, and standout frames
%   \bgpatternblock — used for pblock/palertblock/pexampleblock title bars
%                     (receives the block color as its argument)
\renewcommand{\bgpattern}{\bgpatterntemplate{wcprimary}}
\renewcommand{\bgpatternblock}[1]{\bgpatterntemplate{#1}}
